%%%%%%%%%%%%%%%%%%%%%%%%%%%%%%%%%%%%%%%%%
% University Assignment Title Page 
% LaTeX Template
% Version 1.0 (27/12/12)
%
% This template has been downloaded from:
% http://www.LaTeXTemplates.com
%
% Original author:
% WikiBooks (http://en.wikibooks.org/wiki/LaTeX/Title_Creation)
%
% License:
% CC BY-NC-SA 3.0 (http://creativecommons.org/licenses/by-nc-sa/3.0/)
% 
% Instructions for using this template:
% This title page is capable of being compiled as is. This is not useful for 
% including it in another document. To do this, you have two options: 
%
% 1) Copy/paste everything between \begin{document} and \end{document} 
% starting at \begin{titlepage} and paste this into another LaTeX file where you 
% want your title page.
% OR
% 2) Remove everything outside the \begin{titlepage} and \end{titlepage} and 
% move this file to the same directory as the LaTeX file you wish to add it to. 
% Then add \input{./title_page_1.tex} to your LaTeX file where you want your
% title page.
%
%%%%%%%%%%%%%%%%%%%%%%%%%%%%%%%%%%%%%%%%%

%----------------------------------------------------------------------------------------
%	PACKAGES AND OTHER DOCUMENT CONFIGURATIONS
%----------------------------------------------------------------------------------------

\documentclass[12pt]{article}

\usepackage[utf8]{inputenc}
\usepackage[T1]{fontenc}
\usepackage{lmodern}
\usepackage{parselines}
\usepackage[portuguese]{babel}
\usepackage{graphicx}
\usepackage[document]{ragged2e}
\usepackage{listings}
\usepackage{xcolor}
\usepackage{geometry}
\geometry{
	a4paper,
	total={170mm,257mm},
	left=30mm,
	right=30mm,
	top=30mm,
	bottom=30mm,
}
\usepackage{amsmath}
\usepackage{hyperref}
\usepackage{url}
\usepackage{float}
\usepackage{tabularx}
\usepackage{booktabs}
\usepackage{indentfirst}

\graphicspath{ {images/} }

\colorlet{punct}{red!60!black}
\definecolor{background}{HTML}{EEEEEE}
\definecolor{delim}{RGB}{20,105,176}
\colorlet{numb}{magenta!60!black}


\lstdefinelanguage{json}{
    basicstyle=\normalfont\ttfamily,
    numbers=left,
    numberstyle=\scriptsize,
    stepnumber=1,
    numbersep=8pt,
    showstringspaces=false,
    breaklines=true,
    frame=lines,
    backgroundcolor=\color{background},
    literate=
     *{0}{{{\color{numb}0}}}{1}
      {1}{{{\color{numb}1}}}{1}
      {2}{{{\color{numb}2}}}{1}
      {3}{{{\color{numb}3}}}{1}
      {4}{{{\color{numb}4}}}{1}
      {5}{{{\color{numb}5}}}{1}
      {6}{{{\color{numb}6}}}{1}
      {7}{{{\color{numb}7}}}{1}
      {8}{{{\color{numb}8}}}{1}
      {9}{{{\color{numb}9}}}{1}
      {:}{{{\color{punct}{:}}}}{1}
      {,}{{{\color{punct}{,}}}}{1}
      {\{}{{{\color{delim}{\{}}}}{1}
      {\}}{{{\color{delim}{\}}}}}{1}
      {[}{{{\color{delim}{[}}}}{1}
      {]}{{{\color{delim}{]}}}}{1},
}

\definecolor{dkgreen}{rgb}{0,0.6,0}
\definecolor{gray}{rgb}{0.5,0.5,0.5}
\definecolor{mauve}{rgb}{0.58,0,0.82}

\lstset{frame=tb,
  language=Java,
  aboveskip=3mm,
  belowskip=3mm,
  showstringspaces=false,
  columns=flexible,
  basicstyle={\small\ttfamily},
  numbers=none,
  numberstyle=\tiny\color{gray},
  keywordstyle=\color{blue},
  commentstyle=\color{dkgreen},
  stringstyle=\color{mauve},
  breaklines=true,
  breakatwhitespace=true,
  tabsize=3
}

\hypersetup{
  colorlinks, linkcolor=black
}

\begin{document}

\begin{titlepage}

\newcommand{\HRule}{\rule{\linewidth}{1mm}} % Defines a new command for the horizontal lines, change thickness here

\center % Center everything on the page
 
%----------------------------------------------------------------------------------------
%	HEADING SECTIONS
%----------------------------------------------------------------------------------------

\includegraphics{feup.jpg}

\textsc{\large Agentes e Inteligência Artificial Distribuída}\\[0.8cm] % Major heading such as course name
\textsc{\large 4º ano do Mestrado Integrado em Engenharia Informática e Computação}\\[0.8cm] % Minor heading such as course title

%----------------------------------------------------------------------------------------
%	TITLE SECTION
%----------------------------------------------------------------------------------------

\HRule \\[1.2cm]
{ \huge \bfseries \textit{Simulação de Evacuação com Agentes}}\\[0.6cm] % Title of your document
{ \large \bfseries Relatório Intercalar} \\[0.6cm]
\HRule \\[2cm]
 
%----------------------------------------------------------------------------------------
%	AUTHOR SECTION
%----------------------------------------------------------------------------------------


% If you don't want a supervisor, uncomment the two lines below and remove the section above
\Large \emph{Authors:}\\[0.5cm] \normalsize
Gil \textsc{Domingues}\\[0.1cm]  
- up201304646@fe.up.pt\\[0.1cm]
Pedro \textsc{Pontes}\\[0.1cm]
- up201305367@fe.up.pt\\[2cm] 

%----------------------------------------------------------------------------------------
%	DATE SECTION
%----------------------------------------------------------------------------------------

{\large \today}\\[0cm] % Date, change the \today to a set date if you want to be precise

%----------------------------------------------------------------------------------------
%	TABLE OF CONTENTS & LISTS OF FIGURES AND TABLES
%----------------------------------------------------------------------------------------

\clearpage 

\tableofcontents

\clearpage 
%----------------------------------------------------------------------------------------
%	INTRODUÇÃO
%----------------------------------------------------------------------------------------
\justify\normalsize

\section{Introdução} 

Uma evacuação implica mover pessoas de um dado local devido à ocorrência de uma situação de (potencial) catástrofe. Exemplos incluem a evacuação de um edifício em chamas ou de uma localidade, antes, durante ou após um desastre natural, como uma cheia ou terramoto. Tipicamente, de uma evacuação acabam por resultar feridos ou mesmo mortes, vítimas de espezinhamento. 

Com o aumento da frequência de situações que implicam a evacuação de um elevado número de pessoas num curto espaço de tempo, existe uma consciência acrescida da importância do planeamento dessas situações.

Com efeito, a gestão e organização de multidões em situações de emergência tornou-se uma importante área de estudo ao longo dos últimos anos e desempenha, hoje, um papel importante no desenho de um edifício ou área.

Dados os desafios - quer de ordem prática, quer de ordem financeira - que a realização de simulacros coloca, é cada vez mais comum o uso de técnicas de simulação para estudar estas situações, existindo diversos tipos de sistemas, como as simulações baseadas na dinâmica de fluxo, as simulações baseadas em autómatos e simulações baseadas em agentes.



\newpage
%----------------------------------------------------------------------------------------
%	EUNCIADO 
%----------------------------------------------------------------------------------------

\section{Enunciado}

\subsection{Descrição}

Ocorreu um incêndio, uma inundação, a libertação de um gás nocivo, um qualquer acidente que obriga à evacuação daqueles presentes num dado local. Esse local possui múltiplas saídas de emergência e também obstáculos. Os indivíduos encontram-se distribuídos pelo local, ocupados nas suas tarefas usuais. Aquando da deteção do acidente, todos os indivíduos procuram atingir uma das saídas de emergência, o mais rapidamente possível.

Alguns agentes poderão ser altruístas, no sentido de ajudarem acidentados a deslocarem-se até à saída, outros poderão simplesmente querer "salvar a pele", exibindo um comportamento mais egoísta, conforme se descreve adiante.


\subsection{Objetivos}

Realizado no âmbito da unidade curricular de Agentes e Inteligência Artificial Distribuída, com este trabalho pretende desenvolver-se um programa que permita simular a interação de agentes confinados a um espaço concreto e limitado, podendo o utilizador configurar o local do acidente, especificando, por um lado, o tipo, número e localização dos agentes a evacuar e, por outro, o número e localização de saídas de emergência e obstáculos.


\subsection{Resultados Esperados e Avaliação}

Variando as podendo o utilizador configurar o local do acidente, especificando, por um lado, o tipo, número e localização dos agentes a evacuar e, por outro, o número e localização de saídas de emergência e obstáculos.Adicionalmente, deverá apresentar diversos valores estatísticos, como o tempo médio e máximo de evacuação ou número de feridos.

\section{Ferramentas}
A implementação do programa descrito será realizada Sajas, Repast Symphony…
Para que serve
Descrição das características principais
Realce das funcionalidades relevantes para o trabalho


\section{Especificação}
\subsection{Agentes}
Na implementação do projeto deverá ser usada uma arquitetura de Subsunção, em que os comportamentos são definidos como regras, 

Identificação e caracterização dos diversos tipos de agente a implementar.


\subsection{Interação}
Protocolos de interacção

\subsection{Planeamento}
Para a implementação, importa distinguir a divisão do programa em módulos:
Representação do Espaço;
Representação dos Agentes;
Definição das Interações;

Com base nesta divisão, definiram-se as seguintes etapas:
-Especificação e planeamento;
-Implementação do espaço;
-Implementação dos agentes;
-Implementação das interações entre agentes;
-Exploração de diferentes cenários e recolha e avaliação de métricas.


\section{Conclusão}

Finda esta primeira fase, consideram-se atingidos os objetivos definidos para esta primeira fase: foi feita a descrição do projeto – ...
Cenário e objectivos do trabalho
Estudo da plataforma/ferramenta (de sistemas multi-agente) para desenvolvimento
Agentes, suas estratégias e interacções
Resultados esperados e como avaliá-los

\section{Recursos}
\subsection{Bibliografia}
[1] Almeida, João; Rosseti, Rosaldo; Coelho, António: Crowd Simulation Modeling Applied to Emergency and Evacuation Simulations using Multi-Agent Systems. (2011)

\subsection{Software}

\end{titlepage}
\end{document}